% Team Reference ICPC. Compilar dos veces: pdflatex notebook.tex
\documentclass[8pt,a4paper,landscape]{extarticle}
\usepackage[utf8]{inputenc}
\usepackage[T1]{fontenc}
\usepackage{lmodern}
\usepackage[a4paper,landscape,left=8mm,right=8mm,top=13mm,bottom=8mm,
            headsep=3mm]{geometry}
\usepackage{multicol}
\usepackage{listings}
\usepackage{xcolor}
\usepackage{fancyhdr}
\usepackage{titlesec}
\usepackage{amsmath,amssymb}
\usepackage{enumitem}
\usepackage{etoolbox}

% ---- Cambiar aqui los datos del equipo ----
\newcommand{\equipo}{Nombre del equipo}
\newcommand{\universidad}{Universidad}

\pagestyle{fancy}
\fancyhf{}
\fancyhead[L]{\small\textbf{\equipo} --- \universidad}
\fancyhead[C]{\small Team Reference ICPC}
\fancyhead[R]{\small P\'agina \thepage}
\renewcommand{\headrulewidth}{0.4pt}
\setlength{\headheight}{11pt}

\setlength{\columnsep}{5mm}
\setlength{\columnseprule}{0.2pt}
\setlength{\parindent}{0pt}
\setlength{\parskip}{2pt}
\renewcommand{\contentsname}{\'Indice}

\titleformat{\section}{\large\bfseries}{\thesection.}{0.5em}{}[\vspace{-0.4ex}\hrule]
\titlespacing*{\section}{0pt}{1.2ex}{0.9ex}
\titleformat{\subsection}{\normalsize\bfseries}{\thesubsection}{0.5em}{}
\titlespacing*{\subsection}{0pt}{0.9ex}{0.1ex}
\titleformat{\paragraph}[runin]{\bfseries}{}{0pt}{}[.]
\titlespacing*{\paragraph}{0pt}{0.8ex}{0.5em}

% ---- Colores del codigo (estilo editor claro) ----
\definecolor{colClave}{RGB}{0,0,190}        % palabras clave
\definecolor{colTipo}{RGB}{0,120,130}       % tipos y contenedores STL
\definecolor{colComentario}{RGB}{0,125,0}   % comentarios
\definecolor{colTexto}{RGB}{170,30,30}      % strings y caracteres
\definecolor{colDirectiva}{RGB}{130,0,140}  % #include, #define
\definecolor{colFondo}{RGB}{245,245,245}    % fondo del bloque

\lstset{
  language=C++,
  basicstyle=\ttfamily\fontsize{8.5pt}{9.8pt}\selectfont,
  columns=fixed,
  basewidth=0.5em,
  keepspaces=true,
  breaklines=true,
  showstringspaces=false,
  tabsize=2,
  keywordstyle=\color{colClave}\bfseries,
  keywordstyle={[2]\color{colTipo}},
  morekeywords={[2]{vector,string,pair,array,queue,priority_queue,
    set,map,unordered_map,complex,function,greater,less,
    mt19937,mt19937_64,uniform_int_distribution,__int128,
    SegTree,SCC,Pt,PtD,Circulo,Consulta}},
  commentstyle=\color{colComentario}\itshape,
  stringstyle=\color{colTexto},
  directivestyle=\color{colDirectiva},
  backgroundcolor=\color{colFondo},
  % Sin marco: con frame=single la columna se corta justo despues
  % del borde superior y el titulo queda separado de su codigo.
  xleftmargin=3pt,
  xrightmargin=3pt,
  aboveskip=3pt,
  belowskip=3pt
}
% listings invita a cortar la columna justo antes de cada bloque,
% lo que deja titulos huerfanos; este parche lo prohibe.
\makeatletter
\patchcmd{\lst@Init}{\par\penalty-50\relax}{\par\nobreak}{}
  {\PackageWarning{notebook}{No se pudo parchear listings}}
\makeatother
\newcommand{\codigo}[2]{\subsection{#1}\lstinputlisting{codigo/#2}}
\setlist{nosep,leftmargin=1.2em}

\begin{document}
\begin{multicols*}{2}

{\LARGE\bfseries Team Reference ICPC}\par
{\large \equipo{} --- \universidad}\par
\vspace{1ex}

\textbf{Convenciones.} Todo el c\'odigo es C++17 y asume
\texttt{\#include <bits/stdc++.h>} y \texttt{using namespace std;}
(ver la plantilla). No se usan macros: cada bloque se copia tal cual.
Los \'indices empiezan en 0. Cuando un bloque dice \emph{Requiere},
hay que copiar tambi\'en ese otro bloque. Cada algoritmo fue probado
contra una soluci\'on de fuerza bruta con casos aleatorios
(carpeta \texttt{tests}).

\tableofcontents

\section{Plantilla y utilidades}
\codigo{Plantilla}{varios/plantilla.cpp}
\codigo{N\'umeros aleatorios}{varios/aleatorio.cpp}
\subsection{Stress test (bash)}
\lstinputlisting[language=bash]{codigo/varios/stress.sh}

\section{Estructuras de datos}
\codigo{Segment Tree 4N con lazy}{estructuras/segtree_lazy.cpp}
\codigo{Segment Tree 2N iterativo}{estructuras/segtree_2n.cpp}
\codigo{Fenwick Tree}{estructuras/fenwick.cpp}
\codigo{Sparse Table}{estructuras/sparse_table.cpp}
\codigo{Union-Find (DSU)}{estructuras/dsu.cpp}
\codigo{Sqrt Decomposition}{estructuras/sqrt_decomp.cpp}
\codigo{Algoritmo de Mo}{estructuras/mo.cpp}
\codigo{Heavy-Light Decomposition (y LCA)}{estructuras/hld.cpp}
\codigo{Centroid Decomposition}{estructuras/centroid.cpp}
\codigo{Ordered Set (pb\_ds)}{estructuras/ordered_set.cpp}
\codigo{Li Chao Tree}{estructuras/lichao.cpp}

\section{Grafos}
\codigo{Dijkstra}{grafos/dijkstra.cpp}
\codigo{Biconexas, puentes, articulaci\'on, bridge tree}{grafos/biconexas.cpp}
\codigo{SCC y grafo reducido}{grafos/scc.cpp}
\codigo{2-SAT}{grafos/twosat.cpp}
\codigo{Flujo m\'aximo (Dinic)}{grafos/dinic.cpp}
\codigo{Flujo m\'aximo de costo m\'inimo}{grafos/mincost.cpp}
\codigo{H\'ungaro (asignaci\'on de peso m\'in/m\'ax)}{grafos/hungaro.cpp}
\codigo{Matching bipartito (Hopcroft-Karp)}{grafos/hopcroft_karp.cpp}

\section{Matem\'atica}
\codigo{Aritm\'etica modular}{matematica/aritmetica.cpp}
\codigo{Miller-Rabin}{matematica/miller_rabin.cpp}
\codigo{Pollard Rho}{matematica/pollard_rho.cpp}
\codigo{Euclides extendido, inverso, diof\'antica}{matematica/euclides.cpp}
\codigo{Teorema chino del resto}{matematica/crt.cpp}
\codigo{Baby Step / Giant Step}{matematica/bsgs.cpp}
\codigo{Criba lineal}{matematica/criba.cpp}
\codigo{Gauss (reales y m\'odulo primo)}{matematica/gauss.cpp}
\codigo{FFT}{matematica/fft.cpp}
\codigo{NTT}{matematica/ntt.cpp}
\codigo{Matrices}{matematica/matriz.cpp}
\codigo{Combinatoria}{matematica/combinatoria.cpp}

\section{Strings}
\codigo{Hashing}{strings/hashing.cpp}
\codigo{KMP}{strings/kmp.cpp}
\codigo{Funci\'on Z}{strings/z.cpp}
\codigo{Manacher}{strings/manacher.cpp}
\codigo{Suffix Automaton (DP en el DAG)}{strings/suffix_automaton.cpp}
\codigo{Suffix Array y LCP}{strings/suffix_array.cpp}
\codigo{Aho-Corasick}{strings/aho_corasick.cpp}

\section{Geometr\'ia}
\codigo{Punto y primitivas}{geometria/punto.cpp}
\codigo{Intersecci\'on de segmentos}{geometria/segmentos.cpp}
\codigo{Convex Hull}{geometria/convex_hull.cpp}
\codigo{Puntos antipodales (rotating calipers)}{geometria/antipodal.cpp}
\codigo{C\'irculo m\'inimo}{geometria/circulo_minimo.cpp}
\codigo{Punto en pol\'igono}{geometria/punto_en_poligono.cpp}
\codigo{Par de puntos m\'as cercano}{geometria/par_mas_cercano.cpp}

\section{F\'ormulas y recordatorios}

\paragraph{L\'imites}
\texttt{int} $\approx 2.1\cdot10^9$, \texttt{long long}
$\approx 9.2\cdot10^{18}$, \texttt{\_\_int128} $\approx 1.7\cdot10^{38}$,
\texttt{double} 15--16 d\'igitos. Unas $10^8$ operaciones simples por
segundo. Gu\'ia por tama\~no: $n\le 10$: $O(n!)$; $n\le 20$: $O(2^n n)$;
$n\le 500$: $O(n^3)$; $n\le 5000$: $O(n^2)$; $n\le 10^6$:
$O(n\log n)$; m\'as: $O(n)$.

\paragraph{Primos \'utiles}
$10^9+7$, $10^9+9$, $998244353 = 119\cdot2^{23}+1$, $2^{61}-1$.
M\'aximo de divisores: 240 para $n\le10^6$, 1344 para $n\le10^9$,
103680 para $n\le10^{18}$.

\paragraph{Sumas}
$\sum_{i=1}^n i = \frac{n(n+1)}{2}$,\;
$\sum i^2 = \frac{n(n+1)(2n+1)}{6}$,\;
$\sum i^3 = \left(\frac{n(n+1)}{2}\right)^2$,\;
$\sum_{i=0}^n r^i = \frac{r^{n+1}-1}{r-1}$,\;
$\sum_{i=1}^n \frac1i \approx \ln n + 0.577$
(por eso $\sum_i n/i = O(n\log n)$).

\paragraph{Combinatoria}
\begin{itemize}
\item $\binom{n}{k} = \binom{n-1}{k-1} + \binom{n-1}{k}$;\;
  $\sum_k \binom{n}{k} = 2^n$.
\item Soluciones de $x_1+\dots+x_k = n$ con $x_i \ge 0$:
  $\binom{n+k-1}{k-1}$.
\item Catalan: $C_n = \frac{1}{n+1}\binom{2n}{n}$:
  1, 1, 2, 5, 14, 42, 132, 429.
\item Desarreglos: $D_n = (n-1)(D_{n-1}+D_{n-2})$, $D_0=1$, $D_1=0$.
\item Lucas ($p$ primo): $\binom{n}{k} \equiv \prod \binom{n_i}{k_i}
  \pmod p$ con $n_i, k_i$ los d\'igitos en base $p$.
\item Burnside: \#\'orbitas $= \frac{1}{|G|}\sum_{g\in G}
  |\mathrm{Fix}(g)|$.
\item Inclusi\'on-exclusi\'on: $|\bigcup A_i| = \sum_{S\neq\emptyset}
  (-1)^{|S|+1} |\bigcap_{i\in S} A_i|$.
\end{itemize}

\paragraph{Teor\'ia de n\'umeros}
\begin{itemize}
\item Fermat: $a^{p-1}\equiv 1 \pmod p$. Euler: $a^{\varphi(m)}\equiv 1
  \pmod m$ si $\gcd(a,m)=1$.
\item $\varphi(n) = n\prod_{p\mid n}(1-\frac1p)$;\;
  $\sum_{d\mid n}\varphi(d) = n$.
\item Wilson: $(p-1)! \equiv -1 \pmod p$.
\item $a x + b y = c$ tiene soluci\'on entera si y solo si
  $\gcd(a,b) \mid c$.
\item $\mathrm{lcm}(a,b) = a / \gcd(a,b) \cdot b$.
\end{itemize}

\paragraph{Grafos}
\begin{itemize}
\item Flujo m\'aximo $=$ corte m\'inimo.
\item K\H{o}nig (bipartito): cobertura m\'inima de v\'ertices $=$
  matching m\'aximo; independiente m\'aximo $= V -$ matching.
\item Cubrimiento m\'inimo por caminos disjuntos en un DAG
  $= V -$ matching (grafo bipartito salida/entrada).
\item Dilworth: anticadena m\'axima $=$ m\'inimo n\'umero de cadenas
  que cubren el orden parcial.
\item Hall: hay matching perfecto del lado $A$ si y solo si
  $|N(S)| \ge |S|$ para todo $S \subseteq A$.
\item Euleriano, no dirigido: conexo y todos los grados pares
  (circuito) o exactamente dos impares (camino). Dirigido: grado de
  entrada $=$ grado de salida.
\item Euler (planar conexo): $V - E + F = 2$.
\item Cayley: $n^{n-2}$ \'arboles etiquetados. Kirchhoff: \#\'arboles
  generadores $=$ cualquier cofactor del laplaciano.
\item 2-SAT: $(a \lor b) \equiv (\lnot a \to b) \land (\lnot b \to a)$.
\end{itemize}

\paragraph{Strings}
\begin{itemize}
\item Per\'iodo m\'inimo de $s$: $n - \pi[n-1]$ (es per\'iodo exacto
  si divide a $n$).
\item Substrings distintos: $\frac{n(n+1)}{2} - \sum \mathrm{lcp}$, o
  $\sum_v \mathrm{len}(v) - \mathrm{len}(\mathrm{link}(v))$ en el
  suffix automaton.
\end{itemize}

\paragraph{Geometr\'ia}
\begin{itemize}
\item \'Area de pol\'igono: $\frac12\left|\sum x_i y_{i+1} -
  x_{i+1} y_i\right|$.
\item Pick: $A = I + \frac{B}{2} - 1$; puntos enteros en un segmento:
  $\gcd(|dx|,|dy|) + 1$.
\item Distancia de $p$ a la recta $ab$:
  $\frac{|(b-a)\times(p-a)|}{|b-a|}$.
\item Her\'on: $A = \sqrt{s(s-a)(s-b)(s-c)}$, $s = \frac{a+b+c}{2}$.
  Circunradio $R = \frac{abc}{4A}$, inradio $r = \frac{A}{s}$.
\item Cosenos: $c^2 = a^2+b^2-2ab\cos C$. Senos:
  $\frac{a}{\sin A} = 2R$.
\item Rotar $\theta$: $(x\cos\theta - y\sin\theta,\;
  x\sin\theta + y\cos\theta)$.
\end{itemize}

\paragraph{Antes de enviar}
\begin{itemize}
\item ?`Desborda \texttt{int}? Usar \texttt{long long} en productos y
  sumas.
\item Casos borde: $n = 1$, vac\'io, todos iguales, grafo no conexo.
\item Reiniciar variables globales entre casos de prueba.
\item M\'odulo de negativos: \texttt{((x \% m) + m) \% m}.
\item Recursi\'on profunda ($n \ge 10^5$ en cadena): cuidado con la
  pila.
\item Usar \texttt{'\textbackslash n'} en vez de \texttt{endl}.
\item Comparar reales con tolerancia (\texttt{1e-9}); preferir
  enteros en geometr\'ia.
\end{itemize}

\paragraph{Referencias para ampliar}
KACTL (KTH), Biblioteca de la UFMG, El Vasito (UNC), notebook de
Stanford, cp-algorithms.com.

\end{multicols*}
\end{document}
